\section{Capítulo I: Antecedentes y contexto histórico}
El presente capítulo tiene como propósito reconstruir el escenario que precedió a la creación de la Red Dorsal Nacional de Fibra Óptica (RDNFO), así como explicar el proceso normativo e institucional que dio origen al proyecto. Para ello, se ha organizado el capítulo en tres apartados. En primer lugar, se describe la situación de las telecomunicaciones en el Perú antes de la implementación de la RDNFO, para prestar atención al centralismo de la infraestructura de fibra óptica, al aislamiento tecnológico de las regiones del interior del país y a los altos costos de conectividad que enfrentaban los usuarios. En segundo lugar, se analiza el marco legal y el origen del proyecto, al detenerse en la Ley N° 29904, en el modelo de Asociación Público-Privada elegido para su ejecución y en el proceso que culminó con la adjudicación de la concesión a la empresa Azteca Comunicaciones. Finalmente, se exponen los objetivos iniciales y la visión estratégica que sustentaron el proyecto, referidos a la integración territorial, a la reducción de las tarifas de transporte de datos y a su rol como base para la inclusión digital y la modernización del Estado peruano. Este recorrido permitirá comprender, en capítulos posteriores, por qué las expectativas depositadas en el proyecto no siempre se vieron reflejadas en su ejecución.
\subsection{La situación de las telecomunicaciones en el Perú (Pre-RDNFO)}
Antes de profundizar en cada uno de los problemas específicos que aquejaban a las telecomunicaciones peruanas, conviene situar brevemente el proceso de apertura del mercado que dio origen al escenario descrito en este apartado. Sobre ello, Pacheco (2012) recuerda que en 1994 se vendieron las acciones de la Compañía Peruana de Teléfonos y de ENTEL Perú, al formalizarse la privatización mediante la firma de dos contratos de concesión que fijaron un régimen de tarifas tope por un periodo determinado, al que después se aplicaría un factor de productividad trasladado a favor de los usuarios, junto con determinadas obligaciones de expansión que, en su momento, fueron cumplidas por la empresa ganadora (pp. 33-34). Este antecedente resulta relevante porque explica por qué, más de una década después, la infraestructura de telecomunicaciones persistía en mostrar un patrón de concentración geográfica: el esquema de privatización había resuelto el problema de la escasez de líneas telefónicas en las principales ciudades, pero no generó por sí mismo los incentivos necesarios para extender la red de transporte hacia el interior del país, al dejar pendiente precisamente el problema que años más tarde buscaría resolver la RDNFO.
\subsubsection{El centralismo de la infraestructura de fibra óptica}
Antes de que se concibiera el proyecto de la RDNFO, los estudios técnicos encargados de diagnosticar el estado de las telecomunicaciones en el país advirtieron que el tendido de fibra óptica se concentraba casi exclusivamente en una franja del territorio, al dejar fuera a extensas zonas del interior. Al respecto, Narcizo (2021) señala que el informe elaborado por la Comisión Multisectorial encargada de diseñar el Plan Nacional para el Desarrollo de la Banda Ancha identificó que las redes dorsales de fibra óptica, que constituyen el principal medio de transmisión para la prestación del servicio de banda ancha, se habían desplegado casi en su totalidad en la Costa del país, mientras que en la Sierra dicho despliegue se limitaba a solamente tres departamentos (Cajamarca, Junín, Puno) y en la Selva no existía ninguna red de transporte, lo que representaba una barrera que restringía la masificación de la banda ancha a nivel nacional (p. 12). Este hallazgo evidencia que el centralismo no era una percepción genérica, sino un problema de infraestructura constatado técnicamente, que hacía prever que cualquier proyecto de conectividad debía necesariamente partir por corregir esta desigualdad geográfica antes de plantearse metas de masificación.

La concentración de la infraestructura previa a la intervención estatal tuvo actores con nombres propios, cuyas decisiones comerciales delinearon el mapa de exclusión que el proyecto buscó corregir. Al respecto, Narcizo (2021) detalla que, entre los años 2011 y 2013, la incipiente red de fibra óptica que llegaba solo a setenta y una capitales provinciales era provista principalmente por tres empresas privadas América Móvil, Internexa y, con especial protagonismo, Telefónica, las cuales enfocaron su despliegue casi con exclusividad en el litoral costero y en zonas muy puntuales de la sierra y selva. Esta configuración histórica evidencia claramente la lógica del mercado privado de la época: el accionar de empresas hegemónicas como Telefónica, al priorizar la expansión únicamente en áreas de alta densidad poblacional y rentabilidad económica, demostró el límite natural de la inversión privada y justificó de manera definitiva la necesidad de una intervención estatal para integrar tecnológicamente al resto del territorio.
\subsubsection{El aislamiento tecnológico de las regiones}
El centralismo de la infraestructura tuvo como correlato directo un marcado aislamiento tecnológico de las regiones del interior del país, que se tradujo en una distribución sumamente desigual del acceso a la banda ancha. Sobre este punto, Narcizo (2021) recoge la información contenida en la Resolución Suprema N° 063-2010-PCM, según la cual, a setiembre de 2009, el Perú contaba con 852 900 conexiones de banda ancha a nivel nacional, cifra que en apariencia se encontraba dentro de lo esperado, pero que escondía una alta concentración de dichas conexiones en el departamento de Lima y en la Provincia Constitucional del Callao, que reunían el 64\% del total, seguidos muy de lejos por La Libertad, con apenas 4,96\%, y por Cusco, con 4,33\% (p. 12). Esta desproporción demuestra que buena parte del territorio nacional permanecía prácticamente desconectado de los beneficios asociados a la sociedad de la información, lo que reforzaba la urgencia de una política pública que atendiera de manera prioritaria a las regiones excluidas.

A esta desigualdad se suma una explicación estructural vinculada a la lógica del mercado de telecomunicaciones, que ayuda a entender por qué el aislamiento regional no se corregía por sí solo. Como explica Narcizo (2021), la brecha existente respondía a la distancia entre el mercado actual y potencial del servicio de telecomunicaciones, por un lado, y el denominado mercado de bajo interés para los operadores privados, por otro, en la medida en que el primero ofrecía altos índices de rentabilidad económica mientras que el segundo, pese a tener una alta rentabilidad social, presentaba una rentabilidad económica reducida (p. 5). Esta constatación resulta central para comprender el aislamiento tecnológico de las regiones, pues confirma que dicho fenómeno no era casual, sino consecuencia previsible de un mercado que, dejado a sus propios incentivos, privilegiaría siempre las zonas de mayor densidad poblacional y capacidad de pago en desmedro de las áreas rurales o de preferente interés social.

La necesidad de superar este aislamiento no solo fue advertida por la doctrina y los estudios técnicos previos, sino que pasa a ser la preocupación central del ente rector. Al respecto, el Ministerio de Transportes y Comunicaciones (2023) sostiene que el objetivo prioritario de su gestión actual radica en aprovechar al máximo esta infraestructura para lograr disminuir la desigualdad en el acceso a las telecomunicaciones a nivel nacional. Esta postura confirma que el aislamiento geográfico diagnosticado en los orígenes de la RDNFO exige una intervención estatal constante, en la cual la infraestructura actúa como un igualador social fundamental para el desarrollo tecnológico de las regiones.

\subsubsection{Los altos costos de conectividad}
El tercer elemento que caracterizaba el escenario previo a la RDNFO era el elevado costo del acceso a Internet, derivado en buena medida de la escasa competencia en el segmento de transporte de datos. En relación con ello, Narcizo (2021) señala que uno de los fundamentos que sustentó la implementación del proyecto fue precisamente la expectativa de reducir los altos costos del acceso a Internet, pues antes de su implementación el precio del servicio de transporte ascendía a doscientos dólares americanos por cada Mbps, cifra muy superior a la tarifa que posteriormente ofrecería la propia RDNFO (p. 17). Este dato ilustra con claridad la magnitud del problema que se buscaba resolver, ya que un costo de esa naturaleza resultaba inviable para masificar el servicio entre la población de menores recursos, sobre todo en las zonas rurales identificadas como prioritarias.

Los elevados costos de conectividad no solo obedecían a la falta de competencia, sino también a las condiciones geográficas y demográficas propias del país, que encarecían cualquier despliegue de infraestructura de telecomunicaciones. En ese sentido, Pacheco (2012) advierte que el Perú posee una geografía difícil, una dispersión demográfica que constituye un reto casi imposible de superar para todos los servicios públicos, y una situación similar de precariedad en el suministro de energía eléctrica que debería acompañar el despliegue de redes de telecomunicaciones (p. 36). Este planteamiento permite entender que los altos costos de conectividad no eran un problema exclusivamente regulatorio o comercial, sino que también respondían a factores estructurales del territorio peruano, lo cual explica por qué se consideró necesario un proyecto de infraestructura de alcance nacional financiado con participación del Estado, antes que confiar exclusivamente en la lógica del mercado.

El diagnóstico de estos altos costos se comprende mejor al analizar la composición real de la tarifa. Al respecto, Rosas (2021) detalla que al sumar todos los componentes de la cadena de prestación (salida internacional, redes de transporte, acceso regional y gastos operativos), el precio final ascendía a 287.83 soles mensuales sin IGV por cada Mbps, una cifra considerablemente más alta que el precio de mercado de ese momento (p. 6). Este desglose evidencia de manera concreta cómo los costos se acumulaban a lo largo de toda la cadena técnica, para encarecer drásticamente el servicio que terminaba por pagar el usuario final.

\subsection{Marco legal y origen del proyecto}
El origen del proyecto de la RDNFO no puede entenderse sin hacer referencia al proceso de discusión técnica y política que antecedió a la propia Ley N° 29904. En ese sentido, Pacheco (2012) señala que, luego de varias agendas digitales que no llegaron a concretarse por falta de un compromiso claro de los distintos actores involucrados, el Ministerio de Transportes y Comunicaciones conformó una Comisión encargada de preparar un Plan de Desarrollo de la Banda Ancha durante el periodo 2010-2011, cuyas conclusiones, así como parte de sus integrantes, parecen haber inspirado la formulación del proyecto de ley que dio origen a la norma finalmente aprobada (p. 36). Este dato demuestra que la Ley N° 29904 no fue una decisión repentina ni improvisada, sino el resultado de un proceso de maduración institucional de varios años, aunque, como el propio autor advierte desde su experiencia como integrante de dicha Comisión, ello no garantizó que el mensaje central de ese trabajo técnico se reflejara fielmente en el texto final de la norma.
\subsubsection{La Ley N° 29904 (Ley de promoción de la banda ancha)}
El proyecto de la RDNFO encuentra su origen inmediato en un proceso legislativo que buscó dar respuesta normativa a los problemas de centralismo, aislamiento y altos costos descritos en el apartado anterior. Sobre este proceso, Narcizo (2021) indica que la norma se gestó a partir de la votación de un texto sustitutorio que unificó los Proyectos de Ley N° 688/2011-CR y N° 999/2011-CR, presentados por la Comisión de Transportes y Comunicaciones, y que fue debatido en la Sesión N° 20 del 14 de junio de 2012, ocasión en la que los congresistas sustentaron la necesidad de construir la RDNFO como el gran conducto que permitiría el desarrollo de la banda ancha en el país (p. 13). Este antecedente parlamentario resulta relevante porque muestra que la ley no surgió de manera aislada, sino como resultado de una discusión que integró dos iniciativas legislativas distintas en torno a un mismo objetivo de política pública.

Una vez aprobada, la Ley N° 29904 estableció con claridad el propósito que debía perseguir la política de banda ancha en el país, al fijar el marco dentro del cual se desarrollaría posteriormente la RDNFO. En su artículo 1, la norma dispone que su objeto es impulsar el desarrollo, la utilización y la masificación de la banda ancha en todo el territorio nacional, tanto en la oferta como en la demanda del servicio, al promover el despliegue de infraestructura, servicios, contenidos, aplicaciones y habilidades digitales, como medio que favorezca y facilite la inclusión social, el desarrollo socioeconómico, la competitividad, la seguridad del país y la transformación organizacional hacia una sociedad de la información y el conocimiento (Ley N° 29904, 2012, art. 1). De esta disposición se desprende que la norma no se limitó a declarar la construcción de una infraestructura, sino que buscó articular dicha infraestructura con objetivos más amplios de inclusión y desarrollo, aunque, como se verá en capítulos posteriores, la ejecución del proyecto terminaría por privilegiar la dimensión de infraestructura por encima de estos otros componentes.

\subsubsection{El modelo de Asociación Público-Privada (APP)}
Para materializar el objetivo de la ley, el legislador optó por un esquema de colaboración entre el Estado y el sector privado antes que por la ejecución directa del proyecto por parte de una entidad pública. En efecto, la propia Ley N° 29904 establece en su artículo 8 que el Estado promoverá la inversión y la implementación de la Red Dorsal Nacional de Fibra Óptica, al poder entregarla en concesión mientras mantiene su titularidad, con la finalidad de garantizar el desarrollo económico y la inclusión social, al encargar a la Agencia de Promoción de la Inversión Privada (PROINVERSIÓN) la conducción del proceso de concesión, y al prever que el Estado intervenga de manera subsidiaria únicamente en aquellas zonas donde no participe la inversión privada (Ley N° 29904, 2012, art. 8). Esta disposición confirma que el modelo elegido fue, desde el diseño normativo, una Asociación Público-Privada en la que el Estado se reservaba la titularidad del activo mientras trasladaba a un operador privado las tareas de diseño, financiamiento, despliegue y operación de la red.

El diseño contractual derivado del modelo de Asociación Público-Privada elegido para la RDNFO incorporó, además, una serie de restricciones específicas sobre el tipo de servicio que el concesionario podía prestar, lo cual condicionó su desempeño comercial posterior. Sobre esta cuestión, Rosas (2021), al citar un informe del Banco Mundial, señala que el contrato de concesión suscrito por veinte años establecía que el objeto social del operador únicamente podía incluir la prestación del servicio portador a nivel nacional, sin posibilidad de ofrecer otros servicios ni de atender directamente a los usuarios finales, y que el propio Ministerio de Transportes y Comunicaciones se comprometía a habilitar a dicho operador como el único prestador de dicho servicio (p. 31). Esta restricción contractual resulta especialmente relevante porque revela que el modelo de Asociación Público-Privada no solo definía quién asumiría el riesgo de la inversión, sino que además limitaba de antemano el margen comercial del concesionario, lo cual, como se analizará en capítulos posteriores, terminaría por afectar la sostenibilidad financiera del proyecto.

\subsubsection{La adjudicación a Azteca Comunicaciones}
Tras la publicación de la Ley N° 29904 y de su reglamento, se dio inicio al proceso de selección del operador privado que se encargaría de ejecutar el proyecto bajo el esquema de concesión previsto en la norma. Sobre este proceso, Narcizo (2021) precisa que entre el 31 de mayo y el 1 de junio de 2013 se publicaron los avisos de convocatoria del Concurso de Proyectos Integrales para la entrega en concesión del proyecto de la RDNFO, y que el 23 de diciembre de ese mismo año el Comité de PROINVERSIÓN en Proyectos de Energía e Hidrocarburos adjudicó la Buena Pro del concurso a la empresa Azteca Comunicaciones Perú S.A.C. (p. 17). Este dato permite fijar con precisión el momento en que el proyecto dejó de ser una previsión normativa para convertirse en un compromiso contractual concreto entre el Estado peruano y un operador privado determinado.

La adjudicación de la Buena Pro fue seguida por la suscripción formal del contrato, el cual estableció las obligaciones de inversión, los plazos de la concesión y el alcance territorial a nivel macrorregional. Al integrar la información sobre este hito, se tiene que el 17 de junio de 2014 se suscribió el Contrato de Concesión por veinte años para el diseño y operación de las coberturas Norte, Sur y Centro (Narcizo, 2021); un acuerdo que, según precisa Rosas (2021), demandó una inversión de 333 millones de dólares americanos y permitió que la red entrara oficialmente en operación en mayo de 2015. Esta formalización establece el punto de partida legal y financiero de la relación con Azteca Comunicaciones; un vínculo que, como se analizará en capítulos posteriores, terminaría fracturado por agudas controversias vinculadas al diseño de sus tarifas y a la competencia directa con la infraestructura de otros operadores privados.

\subsection{Objetivos iniciales y visión estratégica}
La visión estratégica que acompañó al proyecto de la RDNFO debe entenderse, asimismo, como parte de una evolución más amplia en la manera en que el Estado peruano concebía el rol de las telecomunicaciones dentro de sus políticas de desarrollo. Al respecto, Narcizo (2021) explica que los proyectos financiados a través del fondo FITEL entre los años 2006 y 2013 transitaron progresivamente desde la instalación de teléfonos públicos hasta la implementación de la banda ancha, tránsito que respondió tanto a la evolución de las tecnologías como a la necesidad de adaptar a la ciudadanía a la nueva era de la información, en línea con el principio de vigencia tecnológica que debe observar todo servicio público (p. 9). Este planteamiento permite comprender que los objetivos de integración territorial, reducción de tarifas y modernización estatal que se desarrollará luego no surgieron de manera aislada, sino que fueron la culminación de una trayectoria de política pública orientada a mantener actualizada la infraestructura de telecomunicaciones del país conforme a las exigencias tecnológicas de cada momento.
\subsubsection{Integración territorial e interconexión}
\subsubsection{Reducción de tarifas de transporte de datos}
\subsubsection{Base para la inclusión digital y modernización del Estado}
